\section{Complete normative Temporale}

The Temporale is generated rather than copied onto fixed civil dates. Let Easter Sunday be day zero. The offsets below locate every movable named observance; the seasonal rules that follow generate every remaining Sunday and feria.

\begin{longtable}{@{}>{\RaggedRight\arraybackslash}p{0.37\linewidth}>{\RaggedRight\arraybackslash}p{0.18\linewidth}>{\RaggedRight\arraybackslash}p{0.12\linewidth}>{\RaggedRight\arraybackslash}p{0.25\linewidth}@{}}
\toprule
\textbf{Observance} & \textbf{Position} & \textbf{Class} & \textbf{Calendar note}\\
\midrule
\endfirsthead
\toprule
\textbf{Observance} & \textbf{Position} & \textbf{Class} & \textbf{Calendar note}\\
\midrule
\endhead
\bottomrule
\endfoot
First Sunday of Advent & fourth Sunday before Christmas & I & Begins the liturgical year.\\
Holy Name of Jesus & Sunday from Jan. 2 through Jan. 5; otherwise Jan. 2 & II & Fixed by a Sunday/date rule.\\
Holy Family & first Sunday after Epiphany & II & Fixed by a Sunday rule.\\
Septuagesima Sunday & Easter $-63$ days & II & Begins Septuagesima season.\\
Sexagesima Sunday & Easter $-56$ days & II & ---\\
Quinquagesima Sunday & Easter $-49$ days & II & ---\\
Ash Wednesday & Easter $-46$ days & I & First weekday of Lent.\\
First through fourth Sundays of Lent & Easter $-42$, $-35$, $-28$, or $-21$ days & I & ---\\
First Sunday of the Passion & Easter $-14$ days & I & Begins Passiontide.\\
Seven Sorrows of the Blessed Virgin Mary & Friday after Passion Sunday & comm. & Seasonal commemoration printed in the calendar.\\
Palm Sunday & Easter $-7$ days & I & Second Sunday of the Passion.\\
Holy Monday, Tuesday, and Wednesday & Easter $-6,-5,-4$ & I & Privileged ferias.\\
Sacred Triduum & Easter $-3,-2,-1$ & I & Holy Thursday, Good Friday, Holy Saturday; highest privileged order.\\
Easter Sunday & Easter & I & With a privileged octave of class I.\\
Monday through Saturday in Easter octave & Easter $+1$ through $+6$ & I & Octave days.\\
Low Sunday & Easter $+7$ & I & Sunday in albis, octave day.\\
Second through fifth Sundays after Easter & Easter $+14$, $+21$, $+28$, or $+35$ days & II & ---\\
Minor Litanies & Monday--Wednesday before Ascension & --- & Classification follows the occurring day; the litany does not create a feast.\\
Office of the Blessed Virgin Mary on Saturday & an otherwise unimpeded Saturday of class IV & IV & A recurring office and Mass, with seasonal forms; not a class-IV feast.\\
Ascension of the Lord & Easter $+39$ & I & Preceded by a class-II vigil.\\
Sunday after Ascension & Easter $+42$ & II & ---\\
Vigil of Pentecost & Easter $+48$ & I & Privileged vigil.\\
Pentecost Sunday & Easter $+49$ & I & With a privileged octave of class I.\\
Pentecost octave, including Ember Days & Easter $+50$ through $+55$ & I & Octave days.\\
Trinity Sunday & Easter $+56$ & I & First Sunday after Pentecost.\\
Corpus Christi & Easter $+60$ & I & Thursday after Trinity Sunday.\\
Sacred Heart of Jesus & Easter $+68$ & I & Friday after the second Sunday after Pentecost.\\
Sundays after Pentecost & weekly through the end of the year & II & The final is always the twenty-fourth Sunday after Pentecost.\\
Christ the King & last Sunday of October & I & Fixed by a Sunday/date rule.\\
\end{longtable}

\subsection{Seasonal generation rules}

\begin{description}[style=multiline,leftmargin=0.25\linewidth,labelwidth=0.22\linewidth,itemsep=0.5em]
\item[Advent] Its four Sundays are class I. Ember Wednesday, Friday, and Saturday and the ferias from December 17 through 23 are class II; the other Advent ferias are class III. December 24 has its own class-I vigil in the fixed table.
\item[Christmas cycle] Christmas has a privileged octave. December 26--31 are class-II octave days, with the saints or commemorations shown in the fixed table; January 1 is the class-I octave day. The Sunday within the octave is class II. Epiphany is class I; the Sundays after Epiphany are class II and ordinary weekday ferias class IV unless another rule supplies a rank.
\item[Septuagesima] Its three Sundays are class II. Weekday ferias are class IV.
\item[Lent and Passiontide] Sundays are class I. Ash Wednesday and Monday through Wednesday of Holy Week are class I. Ember Days are class II. Other Lenten ferias are class III and have the seasonal privileges assigned by the Code. The Sacred Triduum is governed by its own highest-order rules.
\item[Easter cycle] Easter and its octave are class I. Low Sunday is class I; later Sundays through the Sunday after Ascension are class II. The weekday ferias outside the octave are class IV unless another rule intervenes. Ascension and Pentecost are class I; Pentecost has a class-I octave.
\item[After Pentecost] Sundays are class II except the class-I feasts assigned to Sunday. Ordinary weekday ferias are class IV. September Ember Wednesday, Friday, and Saturday are class II.
\end{description}

On an otherwise unimpeded class-IV Saturday, the Office and Mass of the Blessed Virgin Mary on Saturday supply the recurring observance prescribed by General Rubrics nn.~78--79. Its seasonal forms do not turn it into a fourth class of feast.

When Easter is early or late, not all Sundays printed after Epiphany are used in their first position. The omitted post-Epiphany formularies are resumed between the twenty-third and the final Sunday after Pentecost as the Missal directs; the final Sunday remains the twenty-fourth. This device completes the Sunday cycle without pretending that every year has the same number of calendar weeks.

\subsection{Fixed seasonal observances not to overlook}

The Major Litany is observed on April 25, subject to its rubrics when Easter or another privileged day occurs. The Ember Days are the Wednesday, Friday, and Saturday after the third Sunday of Advent, after the first Sunday of Lent, in Pentecost week, and after September 14. These generating rules are part of the calendar even though the Ember Days do not appear as sanctoral feasts.
